\scanpage{222}
\margin{221}
\mspage{L 19}
Récoltes et Semailles j'ai pu me rendre compte, et de façon saisissante en ces tout derniers mois surtout (avec l'écriture des Quatre Opérations), qu'il y a eu depuis mon départ de la scène mathématique une stupéfiante d\enspace é\enspace g\enspace r\enspace a\enspace d\enspace a\enspace t\enspace i\enspace o\enspace n dans l'esprit qui aujourd'hui fait loi dans les milieux \unsure{qye|que} j'avais connus, et (me semble-t-il, dans une large mesure au moins) dans le monde mathématique en général (*). Il est possible même,\deleted{\unsure{que}} tant par ma personnalité mathématique très particulière que par les conditions qui ont entouré mon départ, que celui-ci ait agi comme un catalyseur dans une évolution qui était déjà en train de se faire (**) -- une évolution dont je n'ai\inserted{alors} rien su percevoir (pas plus qu'aucun autre de mes collègues et amis, à la seule exception peut-être de Claude Chevalley). L'aspect de cette dégradation auquel je pense surtout ici (qui\inserted{en} est\inserted{\unsure{juste|partie}} \uline{un} aspect parmi de nombreux autres (***)) est le m\enspace é\enspace p\enspace r\enspace i\enspace s\quad t\enspace a\enspace c\enspace i\enspace t\enspace e, quand ce n'est la dérision sans équivoque, à l'encontre de ce qui (en mathématique, en l'occurrence) ne s'apparente pas au pur travail du marteau sur l'enclume ou sur le burin -- le mépris des processus créateurs les plus délicats (et \deleted{\unsure{parf}} souvent de moindre apparence\inserted{)}; de tout ce qui est \deleted{\unsure{inspiration}} i\enspace n\enspace s\enspace p\enspace i\enspace r\enspace a\enspace t\enspace i\enspace o\enspace n, r\enspace ê\enspace v\enspace e, v\enspace i\enspace s\enspace i\enspace o\enspace n (si puissantes et si fertiles soient-elles), et même (à la limite) de toute i\enspace d\enspace é\enspace e, si clairement conçue et formulée soit-elle : de tout ce qui n'est écrit et p\enspace u\enspace b\enspace l\enspace i\enspace é noir sur blanc, sous forme d'énoncés purs et durs, répertoriables et répertoriés, mûrs pour les ``banques de données'' \deleted{\unsure{s'}}engouffrés dans les inépuisables mémoires de nos mégaordinateurs.

Il y a eu (pour reprendre une expression de C.L. Siegel (\unsure{****}) un extraordinaire \deleted{``\unsure{aplatissement}''}\inserted{``a\enspace p\enspace l\enspace a\enspace t\enspace i\enspace s\enspace s\enspace e\enspace m\enspace e\enspace n\enspace t''}, un ``r\enspace é\enspace t\enspace r\enspace é\enspace c\enspace i\enspace s\enspace s\enspace e\enspace m\enspace e\enspace n\enspace t'' de la pensée mathématique,\deleted{\unsure{f}} dépouillée d'une dimension essentielle, de tout son ``versant d'ombre'', du versant ``féminin''. Il est vrai que par une tradition ancestrale, ce versant-là du travail de découverte restait dans une large mesure\inserted{ }occulté, personne (autant dire) n'en p\enspace a\enspace r\enspace l\enspace a\enspace i\enspace t jamais -- mais le contact\inserted{vivant} avec les sources profondes du rêve, qui alimentent les grandes visions et les grands desseins, n'avait jamais encore (à ma connaissance) été perdu. Il semblerait que\inserted{dès à} présent nous soyons déjà entrés dans une é\enspace p\enspace o\enspace q\enspace u\enspace e\quad d\enspace e\quad d\enspace e\enspace s\enspace s\enspace è\enspace c\enspace h\enspace e\deleted{\unsure{men}}\enspace m\enspace e\enspace n\enspace t, où cette source est, non point tarie certes, mais où l'accès à elle est condamné\deleted{\unsure{s}}, par le verdict sans appel du mépris général et par les représailles de la dérision.

\inserted{\figure{\begin{tikzpicture}[scale=.3]\draw[blue] (0,0)--(1.2,0)--(1.2,1)--(3.3,1);\draw[blue] (1.05,0)--(1.05,1);\end{tikzpicture}}}
Nous voilà approcher du moment, semble-t-il, où sera éradiqué en chacun non seulement le s\enspace o\enspace u\enspace v\enspace e\enspace n\enspace i\enspace r de tout travail proche de la source, du travail ``au féminin'' (ridiculisé
